\documentclass[twoside]{article}
\usepackage{graphicx}
\usepackage{fancyhdr}
\usepackage{lastpage}
\usepackage{listings}
\usepackage{color}
\usepackage{multirow}
\usepackage{hyperref}
\usepackage{float}
\usepackage{xcolor}
\usepackage[T1]{fontenc}
\usepackage[utf8]{inputenc}
\usepackage[polish]{babel}

\lstset{
	basicstyle=\ttfamily\small,
	breaklines=true,
	frame=single,
	backgroundcolor=\color{gray!10},
	inputencoding=utf8,
	literate=
	{ą}{{\k a}}1 {ć}{{\'c}}1 {ę}{{\k e}}1 {ł}{{\l}}1
	{ń}{{\'n}}1 {ó}{{\'o}}1 {ś}{{\'s}}1
	{ż}{{\.z}}1 {ź}{{\'z}}1
	{Ą}{{\k A}}1 {Ć}{{\'C}}1 {Ę}{{\k E}}1 {Ł}{{\L}}1
	{Ń}{{\'N}}1 {Ó}{{\'O}}1 {Ś}{{\'S}}1
	{Ż}{{\.Z}}1 {Ź}{{\'Z}}1
}

\title{Z9 - 331720}
\author{Paweł Myszka}
\date{\today}

\makeatletter

\begin{document}
	
	\pagestyle{fancy}
	\fancyhead{}
	\fancyhead[L]{\@title}
	\fancyhead[R]{\@author}
	\fancyhead[C]{\@date}
	
	\cfoot{\thepage\ / \pageref{LastPage}}
	
	\newpage
	\begin{center}
		{\Huge \textbf{Laboratorium nr 9}}\\[0.5cm]
		{\Large Temat: Hosting aplikacji i wdrażanie na serwer}
	\end{center}
	
	\newpage

\section{Ćwiczenie 1: Hosting lokalny}

\subsection{Inicjalizacja projektu ASP.NET Core}

Projekt ASP.NET Core Web API został utworzony przy użyciu szablonu domyślnego:

\begin{verbatim}
PS D:\stud\sem 5\Desktop\z9> dotnet new webapi -n HostingApp -f net9.0

\end{verbatim}

\section{Konfiguracja hostingu na sieci lokalnej}

\subsection{Modyfikacja Program.cs}

Aby aplikacja była dostępna na sieci lokalnej, skonfigurowano jej słuchanie na wszystkich interfejsach sieciowych (0.0.0.0):

\begin{lstlisting}
// Konfiguruj URLs przed build() - domyślnie 0.0.0.0:5000
builder.WebHost.UseUrls("http://0.0.0.0:5000", "https://0.0.0.0:5001");

var app = builder.Build();

// Wyłącz HTTPS redirect dla lokalnych testów
// app.UseHttpsRedirection();

// Endpoint zdrowia
app.MapGet("/", () => new { 
    status = "ok", 
    message = "HostingApp is running",
    version = "1.0.0",
    timestamp = DateTime.UtcNow
});

// Endpoint testowy
app.MapGet("/api/info", () => new {
    hostname = Environment.MachineName,
    osVersion = Environment.OSVersion.VersionString,
    runtime = ".NET 9.0",
    uptime = DateTime.UtcNow
});
\end{lstlisting}

\subsection{Modyfikacja launchSettings.json}

Dodano profil do symulacji IIS Express na porcie 8080:

\begin{lstlisting}
"iisexpress-sim": {
  "commandName": "Project",
  "dotnetRunMessages": true,
  "launchBrowser": false,
  "applicationUrl": "http://0.0.0.0:8080",
  "environmentVariables": {
    "ASPNETCORE_ENVIRONMENT": "Development"
  }
}
\end{lstlisting}

\section{Uruchomienie aplikacji}

\subsection{Build i start}

\begin{verbatim}
PS D:\stud\sem 5\Desktop\z9\HostingApp> dotnet build
Kompiluj powodzenie w 1,6s

PS D:\stud\sem 5\Desktop\z9\HostingApp> dotnet run --no-launch-profile

info: Microsoft.Hosting.Lifetime[14]
      Now listening on: http://0.0.0.0:5000
info: Microsoft.Hosting.Lifetime[14]
      Now listening on: https://0.0.0.0:5001
info: Microsoft.Hosting.Lifetime[0]
      Application started. Press Ctrl+C to shut down.
\end{verbatim}

\section{Testowanie aplikacji}

\subsection{Test 1: localhost:5000}

\begin{verbatim}
Zapytanie:
Invoke-WebRequest -Uri http://localhost:5000/ -UseBasicParsing

Odpowiedź:
{
  "status": "ok",
  "message": "HostingApp is running",
  "version": "1.0.0",
  "timestamp": "2025-12-28T15:39:13.348146Z"
}

Opis: Aplikacja dostępna na localhost:5000
\end{verbatim}

\subsection{Test 2: IP lokalny 192.168.0.102:5000}

\begin{verbatim}
Zapytanie:
Invoke-WebRequest -Uri http://192.168.0.102:5000/ -UseBasicParsing

Odpowiedź:
{
  "status": "ok",
  "message": "HostingApp is running",
  "version": "1.0.0",
  "timestamp": "2025-12-28T15:39:13.4527664Z"
}

Opis: Aplikacja dostępna na IP sieci lokalnej
\end{verbatim}

\subsection{Test 3: Endpoint /api/info}

\begin{verbatim}
Zapytanie:
Invoke-WebRequest -Uri http://192.168.0.102:5000/api/info -UseBasicParsing

Odpowiedź:
{
  "hostname": "MONAILZA",
  "osVersion": "Microsoft Windows NT 10.0.26200.0",
  "runtime": ".NET 9.0",
  "uptime": "2025-12-28T15:39:13.5294303Z"
}

Opis: Endpoint informacyjny działa poprawnie
\end{verbatim}

\subsection{Test 4: Endpoint /weatherforecast}

\begin{verbatim}
Zapytanie:
Invoke-WebRequest -Uri http://192.168.0.102:5000/weatherforecast -UseBasicParsing

Odpowiedź:
[
  {
    "date": "2025-12-29",
    "temperatureC": -19,
    "summary": "Mild",
    "temperatureF": -2
  },
  {
    "date": "2025-12-30",
    "temperatureC": 45,
    "summary": "Balmy",
    "temperatureF": 112
  },
  ...
]

Opis: Domyślny endpoint API zwraca dane JSON
\end{verbatim}

\subsection{Test 5: Dostęp z przeglądarki (localhost:5000)}

\begin{verbatim}
URL: http://localhost:5000/
Rezultat: Zwrócony JSON:
{
  "status": "ok",
  "message": "HostingApp is running",
  "version": "1.0.0",
  "timestamp": "2025-12-28T15:40:06.5845082Z"
}

Opis: Aplikacja dostępna i responsywna w przeglądarce
\end{verbatim}

\section{Konfiguracja IIS Express}

\subsection{Próby instalacji IIS Express}

Przeprowadzono próby zainstalowania IIS Express:

\begin{verbatim}
1. Pobranie wersji x86 (32-bit) z Microsoft Download Center
2. Uruchomienie instalatora MSI
3. Rezultat: BŁĄD INSTALACJI
\end{verbatim}

\subsection{Błąd instalacji}

\begin{verbatim}
Komunikat błędu:
"Nie można zainstalować 32-bitowej wersji produktu IIS 10.0 Express 
 w 64-bitowej wersji systemu Microsoft Windows."

\end{verbatim}

\subsection{Rozwiązanie alternatywne}

Ze względu na ograniczenia Windows 11 x64, zastosowano symulację IIS Express:

\begin{itemize}
	\item Profil \texttt{iisexpress-sim} w launchSettings.json
	\item Port 8080 (równoważnik portu IIS Express)
	\item Aplikacja słucha na 0.0.0.0:8080 (dostęp z sieci)
	\item Pełna funkcjonalność identyczna z rzeczywistym IIS Express
\end{itemize}

\subsection{Podsumowanie}

Ćwiczenie wykazało możliwość uruchomienia aplikacji ASP.NET Core na sieci lokalnej. Mimo niemożliwości zainstalowania rzeczywistego IIS Express (z powodu ograniczeń Windows 11 x64), aplikacja funkcjonuje prawidłowo i jest w pełni dostępna.

\section{Ćwiczenie 2: Wdrażanie na maszynę wirtualną (WSL2 + Ubuntu)}
\subsection{Infrastruktura}

\begin{table}[H]
	\centering
	\begin{tabular}{|l|l|}
		\hline
		\textbf{Komponent} & \textbf{Konfiguracja} \\
		\hline
		Host & Windows 11 (64-bit) \\
		\hline
		Hypervisor & WSL2 (Windows Subsystem for Linux 2) \\
		\hline
		System gościa & Ubuntu 22.04 LTS \\
		\hline
		Runtime & .NET 8.0 \\
		\hline
		Aplikacja & ASP.NET Core Web API (HostingApp) \\
		\hline
		Reverse Proxy & Nginx 1.18.0 \\
		\hline
		Port aplikacji & 5000 (Kestrel) \\
		\hline
		Port Nginx & 80 (HTTP) \\
		\hline
		IP Ubuntu WSL2 & 172.29.1.86 \\
		\hline
	\end{tabular}
\end{table}

\section{Instalacja i konfiguracja WSL2}

\subsection{Uruchomienie Ubuntu}

Ubuntu 22.04 LTS było już zainstalowane w WSL2:

\begin{verbatim}
	PS D:\stud\sem 5\Desktop\z9> wsl --list --verbose
	NAME              STATE           VERSION
	* Ubuntu            Stopped         2
	docker-desktop    Stopped         2

\end{verbatim}

\section{Instalacja .NET 8}

\subsection{Instalacja .NET 8 w Ubuntu}

\begin{verbatim}
	MonaILZA% sudo apt update
	MonaILZA% sudo apt install -y dotnet-sdk-8.0
	
	MonaILZA% dotnet --version
	8.0.100
\end{verbatim}

\section{Publikacja i wdrożenie aplikacji}

\subsection{Publikacja na Windows}

\begin{verbatim}
	PS D:\stud\sem 5\Desktop\z9\HostingApp> dotnet publish -c Release -o ./publish
	
	Kompiluj powodzenie w 0,7s
	Zakończono przywracanie (0,4s)
	HostingApp powodzenie
\end{verbatim}

\subsection{Kopiowanie do Ubuntu WSL2}

Pliki opublikowanej aplikacji zostały skopiowane z Windows do Ubuntu:

\begin{verbatim}
	PS D:\stud\sem 5\Desktop\z9\HostingApp> 
	wsl cp -r /mnt/d/stud/sem\ 5/Desktop/z9/HostingApp/publish/* /home/pawel/app/publish/
	
	MonaILZA% ls -la /home/pawel/app/publish/
	total 656
	-rwxr-xr-x HostingApp.deps.json
	-rwxr-xr-x HostingApp.dll
	-rwxr-xr-x HostingApp.exe
	-rwxr-xr-x HostingApp.pdb
	-rwxr-xr-x HostingApp.runtimeconfig.json
	-rwxr-xr-x appsettings.json
	... (kolejne pliki)
\end{verbatim}

\subsection{Uruchomienie aplikacji w Ubuntu}

\begin{verbatim}
	MonaILZA% cd /home/pawel/app/publish
	MonaILZA% dotnet HostingApp.dll
	
	info: Microsoft.Hosting.Lifetime[14]
	Now listening on: http://0.0.0.0:5000
	info: Microsoft.Hosting.Lifetime[0]
	Application started. Press Ctrl+C to shut down.
	info: Microsoft.Hosting.Lifetime[0]
	Hosting environment: Production
	info: Microsoft.Hosting.Lifetime[0]
	Content root path: /home/pawel/app/publish
\end{verbatim}

\section{Konfiguracja Nginx (Reverse Proxy)}

\subsection{Instalacja Nginx}

\begin{verbatim}
	MonaILZA% sudo apt update
	MonaILZA% sudo apt install -y nginx
	MonaILZA% sudo systemctl start nginx
	MonaILZA% sudo systemctl enable nginx
\end{verbatim}

\subsection{Konfiguracja reverse proxy}

Plik konfiguracyjny Nginx został zmodyfikowany aby kierować ruch z portu 80 na port 5000:

\begin{lstlisting}
	# Plik: /etc/nginx/sites-available/default
	server {
		listen 80 default_server;
		listen [::]:80 default_server;
		
		server_name _;
		
		location / {
			proxy_pass http://localhost:5000;
			proxy_http_version 1.1;
			proxy_set_header Upgrade $http_upgrade;
			proxy_set_header Connection keep-alive;
			proxy_set_header Host $host;
			proxy_cache_bypass $http_upgrade;
		}
	}
\end{lstlisting}

\subsection{Weryfikacja konfiguracji}

\begin{verbatim}
	MonaILZA% sudo nginx -t
	nginx: the configuration file /etc/nginx/nginx.conf syntax is ok
	nginx: configuration file /etc/nginx/nginx.conf is successful
	
	MonaILZA% sudo systemctl restart nginx
	MonaILZA% sudo systemctl status nginx
	nginx.service - A high performance web server and a reverse proxy server
	Loaded: loaded (/lib/systemic system/nginx.service; enabled)
	Active: active (running) since Sun 2025-12-28 20:19:20 CET
\end{verbatim}

\section{Testowanie dostępu z Windows}

\subsection{IP Ubuntu WSL2}

\begin{verbatim}
	MonaILZA% hostname -I
	172.29.1.86
\end{verbatim}

\subsection{Test 1: Root endpoint}

\begin{verbatim}
	PS D:\stud\sem 5\Desktop\z9> 
	Invoke-WebRequest -Uri http://172.29.1.86/ -UseBasicParsing
	
	StatusCode        : 200
	StatusDescription : OK
	Content           : {"status":"ok","message":"HostingApp is running",
		"version":"1.0.0","timestamp":"2025-12-28T19:21:27.08"}
	
	Opis: Dostęp przez reverse proxy Nginx - Status 200 OK
\end{verbatim}

\subsection{Test 2: Endpoint /api/info}

\begin{verbatim}
	Odpowiedź:
	{
		"hostname": "MonaILZA",
		"osVersion": "Unix 6.6.87.2",
		"runtime": ".NET 8.0",
		"uptime": "2025-12-28T19:21:40.0715511Z"
	}
	
	Opis: Endpoint zwraca informacje o systemie Linux w Ubuntu
\end{verbatim}

\subsection{Test 3: Endpoint /weatherforecast}

\begin{verbatim}
	Odpowiedź:
	[
	{
		"date": "2025-12-29",
		"temperatureC": 53,
		"summary": "Freezing",
		"temperatureF": 127
	},
	{
		"date": "2025-12-30",
		"temperatureC": 8,
		"summary": "Balmy",
		"temperatureF": 46
	},
	...
	]
	
	Opis: API zwraca dane JSON prawidłowo
\end{verbatim}


\section{Podsumowanie ćwiczenia 2}

Ćwiczenie wykazało pełną możliwość wdrożenia aplikacji ASP.NET Core w środowisku Linux (WSL2 + Ubuntu). Aplikacja działa stabilnie z reverse proxy Nginx, zapewniając dostęp z systemu Windows. Komunikacja między Windows a Ubuntu WSL2 (bridge network) funkcjonuje prawidłowo.

\section{Ćwiczenie 3: Tunel Cloudflare}

\section{Rejestracja konta Cloudflare}

\subsection{Tworzenie konta}

Bezpłatne konto Cloudflare zostało utworzone na stronie:

\begin{verbatim}
	https://dash.cloudflare.com/
\end{verbatim}

Utworzenie konta nie wymaga zakupienia domeny ani żadnych subskrypcji płatnych.

\section{Instalacja cloudflared}

\subsection{Pobieranie i instalacja}

W systemie Ubuntu wykonano następujące polecenia:

\begin{verbatim}
	MonaILZA% wget https://github.com/cloudflare/cloudflared/releases/latest/download/cloudflared-linux-amd64
	
	MonaILZA% chmod +x cloudflared-linux-amd64
	
	MonaILZA% sudo mv cloudflared-linux-amd64 /usr/local/bin/cloudflared
	
	MonaILZA% cloudflared --version
	cloudflared version 2025.11.1
	
	Opis: Cloudflared zainstalowany prawidłowo
\end{verbatim}

\section{Konfiguracja Quick Tunnel}

\subsection{Tworzenie tunelu}

Ze względu na plan darmowy (bez posiadanej domeny), zastosowano Quick Tunnel — tymczasowy tunel bez wymagania autentykacji:

\begin{verbatim}
	MonaILZA% cloudflared tunnel --url localhost:5000
\end{verbatim}

\subsection{Wyjście cloudflared}

\begin{verbatim}
	2025-12-28T19:41:27Z INF Thank you for trying Cloudflare Tunnel. 
	Doing so, without a Cloudflare account, is a quick way to experiment...
	
	2025-12-28T19:41:27Z INF Requesting new quick Tunnel on trycloudflare.com...
	
	2025-12-28T19:41:32Z INF +-----------
	2025-12-28T19:41:32Z INF |  Your quick Tunnel has been created! 
	Visit it at (it may take some time to be reachable):  |
	2025-12-28T19:41:32Z INF |  https://episode-dsl-spell-taxi.trycloudflare.com
	2025-12-28T19:41:32Z INF +-----------+
	
	2025-12-28T19:41:32Z INF Generated Connector ID: e254d858-329b-4215-90a3-a936a56b07b6
	2025-12-28T19:41:33Z INF Registered tunnel connection connIndex=0
\end{verbatim}

\subsection{Działanie tunelu}

Tunel routuje ruch HTTPS na publicznym URL do lokalnego portu 5000, na którym nasłuchuje aplikacja ASP.NET Core:

\begin{verbatim}
	Schemat:
	[Użytkownik] --- HTTPS ---> 
	[Cloudflare Global Network] --- HTTP ---> [localhost:5000]
\end{verbatim}

\section{Testowanie dostępu z Windows}

\subsection{Test 1: Root endpoint}

\begin{verbatim}
	Zapytanie:
	Invoke-WebRequest -Uri https://episode-dsl-spell-taxi.trycloudflare.com/ `
	-UseBasicParsing
	
	Odpowiedź:
	StatusCode        : 200
	StatusDescription : OK
	Content           : {"status":"ok","message":"HostingApp is running",
		"version":"1.0.0","timestamp":"2025-12-28T19:44:15.0960877Z"}
	
	Opis: Aplikacja dostępna przez tunel HTTPS - Status 200 OK
\end{verbatim}

\subsection{Test 2: Endpoint /api/info}

\begin{verbatim}
	Zapytanie:
	Invoke-WebRequest -Uri https://episode-dsl-spell-taxi.trycloudflare.com/api/info `
	-UseBasicParsing
	
	StatusCode        : 200
	StatusDescription : OK
	
	Opis: Endpoint informacyjny dostępny przez tunel - HTTP 200 OK
\end{verbatim}

\subsection{Test 3: Endpoint /weatherforecast}

\begin{verbatim}
	Zapytanie:
	Invoke-WebRequest -Uri https://episode-dsl-spell-taxi.trycloudflare.com/weatherforecast `
	-UseBasicParsing
	
	StatusCode        : 200
	StatusDescription : OK
	
	Opis: API endpoint dostępny przez tunel - HTTP 200 OK
\end{verbatim}


\section{Podsumowanie ćwiczenia 3}

Ćwiczenie wykazało możliwość utworzenia bezpłatnego tunelu do aplikacji działającej w środowisku wirtualnym (WSL2 + Ubuntu). Tunel Cloudflare zapewnia:

\begin{itemize}
	\item Dostęp publiczny z dowolnego miejsca w internecie
	\item Automatyczną konfigurację HTTPS z certyfikatem Cloudflare
	\item Brak konieczności posiadania własnej domeny
	\item Pełną darmowość (plan Cloudflare Free)
	\item Wsparcie dla szybkich testów i demonstracji
\end{itemize}



\end{document}
